\subsection{The genericity device (Claim 6.4/6.9)}
\label{sec:molecular-algebraic-induction-genericity}

The accounting nodes above (\cref{lem:case-I}, \cref{lem:case-II},
\cref{thm:theorem-55}, \cref{prop:rigidity-matrix-prop11}) and the
projective assembly of the cycle realization (\cref{lem:cycle-realization})
all reference one shared analytic crux, scoped out into the focused
sub-phase \textbf{Phase~21b} (the analytic sibling of the Phase-21a meet).
The device is the node \cref{lem:genericity-device}, now green.

\medskip\noindent\emph{What the device does and does not discharge.} The
device lifts a rank \emph{attained at one realization} to the same rank at a
generic point. It supplies the genericity inputs each accounting node carries
as a hypothesis (the block-triangular gluing inequality of \cref{lem:case-I},
the span membership of \cref{lem:case-II}, the generic-rank lower bound
$\mathtt{hgen}$ of \cref{prop:rigidity-matrix-prop11}, the projective assembly
of \cref{lem:cycle-realization}). It does \emph{not} build the realizations the
induction consumes. Two things stand between the device and a realization
producer: (i) the device must be \emph{applied} to a varying panel-coordinate
family, which needs the row coordinatization
\cref{lem:rows-polynomial-in-normals} (B0, now green) --- the prior phase only
ever invoked the device on a constant family; and (ii) the device needs a
\emph{seed} realization to lift. The device-to-motive closure itself --- lift a
seed's witnessed independent rows to a generic realization rigid on $V(G)$ --- is
the green glue \cref{lem:realization-of-independent-rows}, shared by both cases;
what remains in each is the \emph{placement}: the seed realization to feed it. For
Case~I the producer (\cref{lem:case-I-realization}) is now green-modulo: KT
eq.~(6.3)'s block-triangular rank-addition over one common framework supplies the
$D(|V(G)|-1)$ independent rows directly, its single residual being KT Claim~6.4
(\cref{lem:claim-6-4}, deferred to Phase~22b); for Case~II
the placement is the $1$-extension normal
(\cref{lem:case-II-realization-placement}), still red. Nor
does the device supply \cref{prop:rigidity-matrix-prop11}'s genericity-free
upper bound $\mathtt{hub}$, a distinct Phase-19-partition obligation carried as
an explicit hypothesis; the device underwrites only the $\mathtt{hgen}$ half.

\begin{lemma}[Genericity device, codimension form; KT Claim 6.4, Claim 6.9]
  \label{lem:genericity-device}
  \lean{CombinatorialRigidity.Molecular.exists_good_realization,
        CombinatorialRigidity.Molecular.exists_good_realization_const}
  \leanok
  \uses{def:rigidity-matrix, def:panel-hinge-framework}
  Fix a multigraph $G$ and regard a panel-hinge realization as a point
  $p \in \mathbb{R}^N$ of the panel-coordinate space (the per-vertex panel
  normals). Each entry of the panel-hinge rigidity matrix $R(G,p)$ is a
  polynomial in $p$ --- of degree two, bilinear in the normals
  (\cref{def:rigidity-matrix}, \cref{def:panel-hinge-framework}). Hence for
  each $r$ the locus $\{p : \operatorname{rank} R(G,p) \ge r\}$, cut out by
  the non-vanishing of some $r \times r$ minor, is Zariski-open, so
  $\operatorname{rank} R(G,p)$ is lower semicontinuous and attains its maximum
  on a dense generic set; a rank attained at one realization is attained
  generically. Dually, in the codimension convention of
  \cref{sec:molecular-rigidity-matrix} ($\operatorname{rank} R(G,p) = D|V| -
  \dim Z(G,p)$, \cref{def:dof-generic}), $\dim Z(G,p)$ is upper semicontinuous
  and minimized generically.
\end{lemma}
\begin{proof}
  \leanok
  A maximal non-vanishing minor of $R(G,p_0)$ is a polynomial in $p$, nonzero
  at $p_0$. Over the infinite field $\mathbb{R}$ a nonzero polynomial does not
  vanish identically (\texttt{MvPolynomial.funext}), so its non-vanishing
  locus is dense and nonempty; on that locus $\operatorname{rank} R(G,p)$ is at
  least the size of the minor. Several generic conditions intersect, each being
  the complement of a proper algebraic subset. See
  \cite[\S6.1, \S6.3]{katohTanigawa2011} (Claim~6.4 and its Case-II twin
  Claim~6.9).

  The formalization is genuinely \emph{multivariate}: the parameter ranges
  over the full panel-coordinate space $\mathbb{R}^N$, not a single affine line.
  The minor's Gram determinant is one multivariate polynomial; nonzero at the
  witness realization, it has a non-root by \texttt{MvPolynomial.funext}, and
  there the rank is at least the witnessed corank. In the codimension shape the
  consumers carry, this is $\#s + \dim Z(F\,p) \le D|V|$ at a produced
  realization. The constant-family closure
  (\texttt{exists\_good\_realization\_const}) reads the bound off a single
  hand-built realization, the form Case~I's block-triangular gluing consumes.
\end{proof}

\noindent The device thus discharges the genericity \emph{hypotheses} the
accounting nodes carry, but not the realization-layer constructions that feed
the induction. With the row coordinatization \cref{lem:rows-polynomial-in-normals}
green and the Case~I producer (\cref{lem:case-I-realization}) now green-modulo KT
Claim~6.4 (\cref{lem:claim-6-4}, deferred to Phase~22b), the still-red
constructions are that residual Claim~6.4 and the Case~II $1$-extension
realization (\cref{lem:case-II-realization}). The device is also consumed by the
Case~III candidate-framework genericity of Phases~22--23.

\subsection{The \texorpdfstring{$V(G)$}{V(G)}-relative count (device output to relative motive)}
\label{sec:molecular-algebraic-induction-relative}

The device (\cref{lem:genericity-device}) produces a generic realization at
which the rigidity-matrix corank is minimized, but its bound is stated
\emph{absolutely}, over the whole screw-assignment space $\R^{D|V_\alpha|}$ of
the ambient body type: a produced point satisfies
$\#s + \dim Z(G,p) \le D|V_\alpha|$. The realization motive
(\cref{def:rank-hypothesis}), by contrast, is $V(G)$-relative
($\mathrm{IsInfinitesimallyRigidOn}\,V(G)$). For a non-spanning inductive
subgraph the two differ by the $D(|V_\alpha| - |V(G)|)$ free dimensions of the
bodies outside $V(G)$ --- each carries no hinge constraint, hence is a free
motion --- so the absolute count alone never reaches the relative motive. The
three nodes here are the adapter the producers (\cref{lem:case-I-realization},
\cref{lem:case-II-realization}) compose with the device to land the relative
motive; they are the ``rank-equality bridge'' \cref{def:rank-hypothesis}
anticipates.

\begin{lemma}[Isolated bodies carry no rigidity: the relative split]
  \label{lem:relative-screw-split}
  \lean{CombinatorialRigidity.Molecular.BodyHingeFramework.finrank_pinnedMotionsOn_vertexSet}
  \leanok
  \uses{def:rigidity-matrix, def:rank-hypothesis, def:pinned-motions-on}
  The screw-assignment space $\R^{D|V_\alpha|} = (\alpha \to \mathrm{Screw})$
  splits as a product along $V(G)$ and its complement, and a body
  $w \in \alpha \setminus V(G)$ carries no hinge constraint, so its $D$ screw
  coordinates are free in every infinitesimal motion. The block pin on the whole
  vertex set $V(G)$ therefore captures exactly the free isolated bodies:
  \[
    \dim\bigl(\mathrm{pinnedMotionsOn}\,V(G)\bigr) \;=\; D\,(|V_\alpha| - |V(G)|),
  \]
  the $\alpha$-dependent free summand of $\dim Z(G,p)$. Combined with the
  block-pin lower bound \cref{lem:pinned-motions-on-rank-bound}, this is what
  isolates the $\alpha$-independent residual rank.
\end{lemma}
\begin{proof}
  \leanok
  An assignment vanishing on $V(G)$ automatically satisfies every hinge
  constraint $S u - S v = 0$ (both endpoints lie in $V(G)$), so
  $\mathrm{pinnedMotionsOn}\,V(G)$ is the kernel of the projection onto the
  $V(G)$ coordinates, $\bigsqcap_{i \in V(G)} \ker(\mathrm{proj}_i)$. Mathlib's
  \texttt{iInfKerProjEquiv} identifies that kernel with the product over the
  complement $\alpha \setminus V(G)$, of dimension $D(|V_\alpha|-|V(G)|)$ by
  \texttt{finrank\_pi\_const}.
\end{proof}

\begin{lemma}[The genericity device, $V(G)$-relative form]
  \label{lem:relative-device-count}
  \lean{CombinatorialRigidity.Molecular.PanelHingeFramework.exists_relative_full_count_ofParam}
  \leanok
  \uses{lem:genericity-device, lem:relative-screw-split,
        lem:rows-polynomial-in-normals}
  Substituting the relative split \cref{lem:relative-screw-split} for the
  absolute screw-count in the device \cref{lem:genericity-device} (applied to
  the varying panel family via the row coordinatization
  \cref{lem:rows-polynomial-in-normals}, B0), a generic panel assignment with at
  least $\#s \ge D(|V(G)|-1)$ witnessed independent rows attains
  \[
    \dim Z(G,p) \;\le\; D\,\bigl(|\alpha \setminus V(G)| + 1\bigr),
  \]
  the device's absolute codimension bound with the $D(|V(G)|-1)$ rows cancelled
  --- i.e.\ the relative full count, the ambient $D|\alpha \setminus V(G)|$ free
  dimensions plus the $D$ trivial-motion dimensions and no residual relative
  corank. Mechanical given \cref{lem:relative-screw-split}: re-wrap the device's
  conclusion using $|\alpha| = |V(G)| + |\alpha \setminus V(G)|$.
\end{lemma}
\begin{proof}
  \leanok
  The device (\cref{lem:rows-polynomial-in-normals}) gives
  $\#s + \dim Z(G,p) \le D|\alpha|$ at the produced point; substitute
  $|\alpha| = |V(G)| + |\alpha \setminus V(G)|$ and the count
  $\#s \ge D(|V(G)|-1)$, then cancel.
\end{proof}

\begin{lemma}[Relative count $\Rightarrow$ the relative motive]
  \label{lem:isInfRigidOn-of-relative-count}
  \lean{CombinatorialRigidity.Molecular.BodyHingeFramework.isInfinitesimallyRigidOn_vertexSet_of_finrank_le}
  \leanok
  \uses{lem:relative-device-count, lem:relative-screw-split,
        lem:pinned-motions-on-rank-bound, lem:rank-delete-vertex, lem:case-I}
  A realization at the relative full count ($k'=0$ in
  \cref{lem:relative-device-count}, i.e.\ $\dim Z(G,p) \le
  D(|\alpha \setminus V(G)| + 1)$) is infinitesimally rigid on $V(G)$:
  $\mathrm{IsInfinitesimallyRigidOn}\,V(G)$. This is the adapter the producers
  compose with the device output, routed through the relative split
  \cref{lem:relative-screw-split} and the relativized Case~I bridge
  \cref{lem:case-I}
  ($\mathrm{isInfinitesimallyRigidOn\_iff\_pinnedMotionsOn\_le}$).
\end{lemma}
\begin{proof}
  \leanok
  Pick any $v_0 \in V(G)$ and read rigidity-on-$V(G)$ off the Case~I bridge
  \cref{lem:case-I} at the trivially-rigid singleton block $\{v_0\}$: it reduces
  to $\mathrm{pinnedMotionsOn}\,\{v_0\} \le \mathrm{pinnedMotionsOn}\,V(G)$. The
  reverse containment is automatic, so it suffices to match dimensions. The
  single-body pin has $\dim(\mathrm{pinnedMotions}\,v_0) = \dim Z(G,p) - D$
  (\cref{lem:rank-delete-vertex}), capped at $D|\alpha \setminus V(G)|$ by the
  hypothesis; the block pin on $V(G)$ has exactly that dimension
  (\cref{lem:relative-screw-split}). Equal dimensions on a containment force the
  needed equality.
\end{proof}

\begin{lemma}[Realization glue from a witnessed independent row family]
  \label{lem:realization-of-independent-rows}
  \lean{CombinatorialRigidity.Molecular.PanelHingeFramework.hasFullRankRealization_of_independent_panelRow}
  \leanok
  \uses{def:rank-hypothesis, lem:relative-device-count,
        lem:isInfRigidOn-of-relative-count, lem:rows-polynomial-in-normals}
  The honest device closure both case producers reduce to once their geometry is
  placed. Fix a multigraph $G$ on a nonempty body set $V(G)$ and an endpoint
  selector $\mathrm{ends}$ recording each edge's link. If at \emph{one} free
  normal assignment $q_0$ the rigidity rows of the free-normal panel family
  $\mathrm{ofNormals}\,G\,\mathrm{ends}\,q_0$ indexed by $s$ are linearly
  independent and $\#s \ge D(|V(G)|-1)$, then $G$ has a full-rank panel
  realization ($\mathrm{HasFullRankRealization}$, the $V(G)$-relative rank).
  This is the composition $\text{N2} \circ \text{N3}$: the device's relative
  count \cref{lem:relative-device-count} lifts the witnessed corank at $q_0$ to a
  generic normal assignment $q$, and the relative-count adapter
  \cref{lem:isInfRigidOn-of-relative-count} turns that into rigidity on
  $V(G) = V(\mathrm{ofNormals}\,G\,\mathrm{ends}\,q)$. The witnessed independent
  family $(q_0, s)$ is the \emph{satisfiable} geometric input the placement
  supplies, not the rank the lemma concludes --- the honesty split matching the
  Case~I splice glue \cref{lem:case-I-splice-seed}.
\end{lemma}
\begin{proof}
  \leanok
  Apply \cref{lem:relative-device-count} to the witness $(q_0, s)$ to obtain a
  normal assignment $q$ with $\dim Z(G,q) \le D(|\alpha \setminus V(G)| + 1)$,
  then \cref{lem:isInfRigidOn-of-relative-count} on
  $\mathrm{ofNormals}\,G\,\mathrm{ends}\,q$ (whose vertex set is $V(G)$).
\end{proof}

\begin{lemma}[Case II placement: the $1$-extension normal; KT (6.12)]
  \label{lem:case-II-realization-placement}
  \uses{def:panel-hinge-framework, lem:case-II, lem:genericity-device,
        lem:exists-independent-panel-extensor,
        lem:case-II-placement-new-rows, lem:case-II-placement-old-rows,
        lem:case-II-placement-old-rows-extract,
        lem:case-II-placement-block-independent,
        lem:case-II-placement-e0-recovery}
  The remaining geometric obligation of \cref{lem:case-II-realization}: from the
  inductive realization of the splitting-off $G_v^{ab}$, construct a free normal
  assignment $q_0$ for the parent graph $G$ --- in particular a panel normal $n$
  for the re-inserted body $v$ --- together with a linearly independent family of
  $D(|V(G)|-1)$ rigidity rows of $\mathrm{ofNormals}\,G\,\mathrm{ends}\,q_0$.
  Geometrically: re-insert $v$ joined to its two neighbours $a, b$ by hinges
  whose supporting extensors are in general position
  (\cref{lem:exists-independent-panel-extensor}), so the two new edges contribute
  $+(D-1)$ independent rows that, combined with the inductive realization's rows
  threaded through the common subgraph $G - v$ (recovering the short-circuit edge
  $e_0$'s constraint from $v$'s two new edges --- the edge substitution
  $G_v^{ab}$ adds $e_0$ but deletes $v$'s edges, so the count is \emph{not} a free
  transport), witness the full target count. This is the shallower parallel of
  the Case~I splice placement (one re-inserted body, $+(D-1)$ rows). \emph{Red}
  (Phase~21b realization layer): the witnessed-independent-row construction across
  the edge substitution is the genuine content; once green it feeds the glue
  \cref{lem:realization-of-independent-rows}.

  \medskip\noindent\emph{Decomposition (Phase-21b realization-layer plan, revised
  2026-06-04).} Per the hand-off discipline for a genuinely-geometric node, this
  placement breaks into buildable sub-nodes plus a closing assembly, mirroring the
  Case~I splice placement's boundary-panel/block-triangular split. The seed family
  $s$ partitions into the \emph{new-edge block} (the two hinges re-attaching $v$)
  and the \emph{old block} (the inductive realization's rows, read off the common
  subgraph $G - v$); the count is $(D-1) + D(|V(G)|-2) = D(|V(G)|-1)$.
  \begin{enumerate}
    \item \cref{lem:case-II-placement-new-rows} (N7b-1, green): choose the normal $n$
      for $v$ so that $v$'s two new hinges supply $D-1$ linearly independent panel rows.
    \item \cref{lem:case-II-placement-old-rows-extract} (N7b-0, green): \emph{produce}
      the old block --- extract $D(|V(G)|-2)$ independent panel rows from the inductive
      realization's rigidity on $V(G_v^{ab})$ (the rank-equals-codimension forward
      direction, the converse of \cref{lem:isInfRigidOn-of-relative-count}).
    \item \cref{lem:case-II-placement-old-rows} (N7b-2, green): transport the
      $e_0$-\emph{free} subfamily of that block onto $G$ through the common subgraph
      $G - v$.
    \item \cref{lem:case-II-placement-block-independent} (N7b-3, green): the two blocks
      are jointly independent (the new-edge rows live in the $v$-column block, the
      old rows in its complement; pin-a-body Lemma~5.1,
      \cref{lem:rank-delete-vertex}, makes the union block-triangular).
    \item \cref{lem:case-II-placement-e0-recovery} (N7b-4, \emph{superseded}): a
      scrutiny pass found N7b-0/1/2/3 do not compose --- N7b-0 extracts an old block
      with no control over $e_0$-rows, but N7b-2 transports only the $e_0$-free rows,
      so the count can fail --- and a recon pass then found the row-side fix
      (``produce an $e_0$-free old block of full count'') \emph{geometrically
      unbuildable} ($G - v$ is not rigid, so no such block exists in $G_v^{ab}$). The
      genuine rank-lift is re-framed \emph{motion-side} (M1--M3,
      \cref{lem:case-II-placement-motion-side-assembly}), which supersedes this whole
      row-side assembly path for Case~II. See ``The rank-lift, motion-side'' below.
  \end{enumerate}
  \emph{Superseded (2026-06-04):} the row-side assembly below was the original plan,
  but the missing producer N7b-4 is geometrically unbuildable in row-side form; the
  Case~II producer instead routes through the motion-side re-frame
  (\cref{lem:case-II-placement-motion-side-assembly}, M3), which concludes rigidity
  on $V(G)$ directly without a witnessed independent row family. The row-side
  sub-nodes N7b-0/1/2/3 + glue N7a stay green and are reused by Case~I.
\end{lemma}
\begin{proof}
  \uses{lem:case-II, lem:exists-independent-panel-extensor,
        lem:case-II-placement-new-rows, lem:case-II-placement-old-rows,
        lem:case-II-placement-old-rows-extract,
        lem:case-II-placement-block-independent,
        lem:case-II-placement-e0-recovery}
  Take $q_0$ to agree with the inductive realization's normals on
  $V(G)\setminus\{v\} = V(G_v^{ab})$ and to assign $v$ the general-position normal
  $n$ of \cref{lem:case-II-placement-new-rows}. Let $s$ be the union of the
  new-edge block (\cref{lem:case-II-placement-new-rows}, $D-1$ rows) and the
  $e_0$-free old block (\cref{lem:case-II-placement-e0-recovery,
  lem:case-II-placement-old-rows}, $D(|V(G)|-2)$ rows);
  \cref{lem:case-II-placement-block-independent} makes the union independent, and
  the count is $D(|V(G)|-1)$. The genuine content is the edge substitution: $s$ is
  \emph{not} a free transport of $G_v^{ab}$'s rows (the short-circuit $e_0$ is
  deleted in $G$, $v$'s two edges added), so the old block is recovered through the
  common subgraph $G - v$ ($\mathrm{removeVertex}$), not along an inclusion
  $G_v^{ab}\le G$; the $e_0$-row's constraint comes from $v$'s two new edges
  (\cref{lem:case-II-placement-e0-recovery}). Deferred to the Phase-21b realization
  layer; see \cite[\S6.3]{katohTanigawa2011} (Lemma~6.8).
\end{proof}

\begin{lemma}[N7b-1: a re-inserted body's transversal hinge gives $D-1$ independent panel rows]
  \label{lem:case-II-placement-new-rows}
  \lean{CombinatorialRigidity.Molecular.BodyHingeFramework.exists_independent_panelRow_of_edge,
        CombinatorialRigidity.Molecular.BodyHingeFramework.exists_independent_panelRow_subfamily_of_edge}
  \leanok
  \uses{def:panel-hinge-framework, def:hinge-row-block, def:rigidity-matrix, lem:rank-delete-vertex}
  Re-inserting the degree-2 body $v$ joins it to its neighbours $a, b$ by two hinges. Each such
  hinge $e = uv$ that is \emph{transversal} --- distinct endpoints and a nonzero supporting extensor
  $C(p(e))$, the general-position condition supplied by choosing $v$'s normal $n$ off the
  support-extensor locus (\cref{lem:exists-independent-panel-extensor}) --- contributes a linearly
  independent family of $D-1 = \mathrm{screwDim}\,k - 1$ rigidity rows of
  $\mathrm{ofNormals}\,G\,\mathrm{ends}\,q_0$, each a member of \emph{that single edge's} panel-row
  span $\langle \mathrm{panelRow}\,\mathrm{ends}\,(e, \cdot, \cdot)\rangle$. These are the $+(D-1)$
  rows the $1$-extension adds in $v$'s column screw block: the hinge-row block $r(p(e))$ is
  $(D-1)$-dimensional (\cref{lem:rank-delete-vertex}, the rank-count input), and its basis lifts
  through the relative-screw evaluation $\mathrm{hingeRow}\,u\,v$ into the per-edge panel-row span.
  This is the panel-row form of the per-edge brick of \cref{def:rigidity-matrix}
  ($\mathrm{exists\_independent\_rigidityRows\_of\_edge}$), restricted to
  membership in this edge's panel rows so the placement assembly
  (\cref{lem:case-II-realization-placement}) can thread it into the device-consuming
  $\mathrm{panelRow}$ family of \cref{lem:realization-of-independent-rows}.

  \medskip\noindent\emph{Honest index-subfamily form
  ($\mathrm{exists\_independent\_panelRow\_subfamily\_of\_edge}$).} The glue
  \cref{lem:realization-of-independent-rows} consumes not a family of \emph{members} of the per-edge
  span but a literal $\mathrm{panelRow}\,\mathrm{ends}$-subfamily indexed by a \emph{set} of panel-row
  indices. This companion form bridges that gap: a transversal hinge $e$ yields an index subset
  $s \subseteq \{e\}\times(\Lambda^k\text{-pairs})$ of size $D-1$ whose \emph{actual}
  $\mathrm{panelRow}\,\mathrm{ends}$-subfamily is independent --- the honest input
  \cref{lem:case-II-realization-placement} threads into the genericity device.
\end{lemma}
\begin{proof}
  \leanok
  \uses{def:hinge-row-block, lem:rank-delete-vertex}
  The transversal hinge's row block $r(p(e)) = (\mathrm{span}\,C(p(e)))^\perp$ has dimension $D-1$
  (\cref{lem:rank-delete-vertex}); pick a basis of it. Its image under the relative-screw evaluation
  $\mathrm{hingeRow}\,u\,v = (\mathrm{screwDiff}\,u\,v)^\ast$ stays independent (the dual map is
  injective at distinct endpoints) and, by the per-edge span identity
  $\mathrm{span\_panelRow\_edge\_eq}$ (the annihilator-family span
  $\mathrm{span\_annihRow\_eq\_dualAnnihilator}$ carried through $\mathrm{hingeRow}\,u\,v$), each row
  lies in the single edge's panel-row span. For the index-subfamily form, the per-edge family
  $(t_1,t_2)\mapsto\mathrm{panelRow}\,\mathrm{ends}\,(e,t_1,t_2)$ spans a $(D-1)$-dimensional space, so
  extracting an independent subfamily of its generators at the span's $\mathrm{finrank}$
  ($\mathrm{Submodule.exists\_fun\_fin\_finrank\_span\_eq}$) and re-indexing each chosen row by its
  $\Lambda^k$-pair gives the honest index subset. See \cite[\S6.3]{katohTanigawa2011}.
\end{proof}

\begin{lemma}[N7b-2: the inductive rows transport through the common subgraph $G-v$]
  \label{lem:case-II-placement-old-rows}
  \lean{CombinatorialRigidity.Molecular.PanelHingeFramework.exists_independent_panelRow_transport}
  \leanok
  \uses{def:panel-hinge-framework, lem:case-II, lem:rows-polynomial-in-normals,
        lem:case-II-placement-old-rows-extract}
  Given the inductive realization's $D(|V(G)|-2)$ independent rows of
  $\mathrm{ofNormals}\,G_v^{ab}\,\mathrm{ends}'\,q_0'$ (produced by
  \cref{lem:case-II-placement-old-rows-extract} from the inductive rigidity),
  transport them to $D(|V(G)|-2)$ independent rows of
  $\mathrm{ofNormals}\,G\,\mathrm{ends}\,q_0$ on $G$. Concretely: along an \emph{injective} reindex $f : s_2 \to s_1$ selecting the
  $e_0$-free subfamily, with each selected row matching across the graph swap (the per-row
  hypothesis $\mathrm{hrow}$), the family is again linearly independent as rows of
  $\mathrm{ofNormals}\,G\,\mathrm{ends}\,q_0$. The transport is \emph{not} along an
  inclusion (neither $G_v^{ab}$ nor $G$ is a subgraph of the other): both sit above the
  common subgraph $G - v$ ($\mathrm{removeVertex}\,v$; $\mathrm{removeVertex\_le}$ and
  $\mathrm{removeVertex\_le\_splitOff}$, green), and the inductive rows that avoid $v$
  --- equivalently avoid the short-circuit edge $e_0$ --- are exactly the rows of
  $G - v$, which survive into $G$ unchanged. The per-row match is where the common
  subgraph enters: a panel row of $\mathrm{ofNormals}$ reads only $\mathrm{ends}$ and the
  normal assignment, not the underlying graph
  (\cref{lem:rows-polynomial-in-normals}, via $\mathrm{toBodyHinge\_supportExtensor}$), so
  when the assembly (\cref{lem:case-II-realization-placement}) picks $q_0$ extending the
  inductive normals and $\mathrm{ends}$ agreeing on $G - v$'s edges, the match is
  definitional. The short-circuit edge $e_0$'s constraint is \emph{not} transported here;
  it is recovered from $v$'s two new edges in N7b-1.
\end{lemma}
\begin{proof}
  \leanok
  \uses{lem:case-II, lem:rows-polynomial-in-normals}
  Restrict the inductive independent family to the rows indexed by edges of
  $G - v$ (drop the $e_0$-row), and read them as rows of $G$ through
  $\mathrm{removeVertex\_le}$. Independence is inherited as a subfamily of an
  independent family ($\mathrm{LinearIndependent.comp}$ along the injective reindex,
  the family equality given by the per-row match). See
  \cite[\S6.3]{katohTanigawa2011} (Lemma~6.8).
\end{proof}

\begin{lemma}[N7b-0: a rigid realization carries a full-rank independent row family]
  \label{lem:case-II-placement-old-rows-extract}
  \lean{CombinatorialRigidity.Molecular.BodyHingeFramework.exists_independent_panelRow_subfamily_of_rigidOn,
        CombinatorialRigidity.Molecular.BodyHingeFramework.finrank_infinitesimalMotions_of_isInfinitesimallyRigidOn_vertexSet}
  \leanok
  \uses{def:panel-hinge-framework, def:rigidity-matrix, def:rank-hypothesis,
        lem:rank-delete-vertex, lem:relative-screw-split}
  The \emph{producer} of the old block that the transport
  \cref{lem:case-II-placement-old-rows} consumes: from the inductive realization
  of $G_v^{ab}$ --- a panel-hinge framework infinitesimally rigid on its own
  vertex set $V(G_v^{ab}) = V(G)\setminus\{v\}$ ($\mathrm{HasFullRankRealization}$,
  the realization motive) --- extract a linearly independent family of
  $D(|V(G_v^{ab})|-1) = D(|V(G)|-2)$ \emph{actual} $\mathrm{panelRow}$ rows of
  that realization (in the honest index-subfamily form
  \cref{lem:case-II-placement-new-rows} supplies for the new block). This is the
  \emph{forward converse} of the relative-count adapter
  \cref{lem:isInfRigidOn-of-relative-count}: that node turns a witnessed corank
  $\#s \ge D(|V|-1)$ \emph{into} rigidity on $V(G)$; this node runs the
  implication backward --- rigidity on $V(G_v^{ab})$ forces
  $\dim Z = D(|V_\alpha| - |V(G_v^{ab})| + 1)$
  (\cref{lem:relative-screw-split}), so $\mathrm{rank}\,R = D|V_\alpha| - \dim Z =
  D(|V(G_v^{ab})|-1)$, and the rigidity rows, spanning a space of that dimension,
  contain an independent subfamily of size $D(|V(G_v^{ab})|-1)$
  ($\mathrm{Submodule.exists\_fun\_fin\_finrank\_span\_eq}$ on the rigidity-row
  span, re-indexed into a $\mathrm{panelRow}$-index subset as in
  \cref{lem:case-II-placement-new-rows}'s honest form). This
  rank-equals-codimension forward direction is the genuine deficit between
  ``rigid on a set'' and ``carries that many independent rows'' --- it is
  \emph{not} discharged by the transport \cref{lem:case-II-placement-old-rows}
  (which only carries an \emph{already-witnessed} independent family across the
  graph swap) nor by the graph-forest existence \cref{def:rigidity-matrix}
  ($\mathrm{exists\_independent\_rigidityRows\_of\_forest}$, which sources rows
  from a transversal-hinge spanning forest, of count bounded by the forest size
  rather than the full $D(|V|-1)$ unless the block is itself rigid).
\end{lemma}
\begin{proof}
  \leanok
  \uses{lem:rank-delete-vertex, lem:relative-screw-split}
  The realization is rigid on $V(G_v^{ab})$, so $\mathrm{pinnedMotionsOn}\,
  V(G_v^{ab})$ pins all residual freedom: its dimension is
  $D(|V_\alpha| - |V(G_v^{ab})|)$ (\cref{lem:relative-screw-split}), and the
  single-body pin gives $\dim Z = D + D(|V_\alpha| - |V(G_v^{ab})|)$
  (\cref{lem:rank-delete-vertex}), whence $\mathrm{rank}\,R = D|V_\alpha| - \dim Z
  = D(|V(G_v^{ab})|-1)$. The rigidity-row span has exactly that dimension, so
  $\mathrm{Submodule.exists\_fun\_fin\_finrank\_span\_eq}$ extracts an independent
  subfamily of that many actual rows; the $\mathrm{panelRow}$-index re-indexing is
  as in \cref{lem:case-II-placement-new-rows}'s honest form. See
  \cite[\S6.3]{katohTanigawa2011} (Lemma~6.8).
\end{proof}

\begin{lemma}[N7b-3: the new-edge and old blocks are jointly independent (pin-a-body column split)]
  \label{lem:case-II-placement-block-independent}
  \lean{CombinatorialRigidity.Molecular.BodyHingeFramework.linearIndependent_sum_pinned_block}
  \leanok
  \uses{def:panel-hinge-framework, lem:rank-delete-vertex}
  The pin-a-body column decomposition (Lemma~5.1, \cref{lem:rank-delete-vertex})
  certifying that the new-edge block (\cref{lem:case-II-placement-new-rows}, the
  rows carried by the re-inserted body $v$'s incident hinges, in the $v$-column
  screw block) and the old block (\cref{lem:case-II-placement-old-rows}, the rows
  of the common subgraph $G-v$, supported off $v$'s columns) form a jointly
  linearly independent family. The split is the body-$v$ screw column: a screw
  assignment pinned to $v$ alone ($\mathrm{Function.update}\,0\,v\,x$) reads the
  old rows as $0$ (their edges avoid $v$) while reading the new rows through $v$'s
  screw column $\mathrm{rn}\,i \circ \mathrm{single}_v$. So if the new rows are
  independent \emph{as functionals of $v$'s screw} (the $D-1$ column-block rows of
  N7b-1) and the old rows are independent (the $D(|V(G)|-2)$ inductive rows of
  N7b-2), their union $\mathrm{Sum.elim}\,\mathrm{rn}\,\mathrm{ro}$ is independent.
  This is the panel-row form of the $1$-extension's exact $+D$ rank lift the
  Case~II accounting \cref{lem:case-II} pins; the assembly
  \cref{lem:case-II-realization-placement} (N7b) supplies the two blocks (whence
  the count $D(|V(G)|-1)$) and discharges the column-split hypotheses by
  construction.
\end{lemma}
\begin{proof}
  \leanok
  \uses{lem:rank-delete-vertex}
  A vanishing combination of the union, evaluated at the body-$v$ pin
  $\mathrm{Function.update}\,0\,v\,x$, collapses to the new block alone (the old
  terms vanish, edges $e_0$-free away from $v$), forcing the new coefficients to
  vanish by their pinned independence; the residual is then a vanishing
  combination of the old block, forcing the old coefficients by their
  independence (pin-a-body Lemma~5.1, \cref{lem:rank-delete-vertex}). See
  \cite[\S6.3]{katohTanigawa2011}.
\end{proof}

\subsubsection*{The rank-lift, motion-side (the N7b-4 re-frame, 2026-06-04)}

A scrutiny+recon pass on the N7b assembly (Phase 21b, 2026-06-04) found the
\emph{row-side} framing of the rank-lift below
(\cref{lem:case-II-placement-e0-recovery}, ``produce an $e_0$-free old block of
full count'') is \emph{geometrically wrong}, and re-frames the genuine content
\emph{motion-side}. The finding, the corrected statement, and the buildable
sub-node decomposition follow; the row-side node is kept (struck) for the
audit trail.

\medskip\noindent\emph{Why the row-side framing fails.} An $e_0$-free old block
of $G_v^{ab}$ is a family of rigidity rows indexed by edges of $G_v^{ab}$ other
than $e_0$ --- i.e.\ rows of the common subgraph $G - v$. But $G - v$ is
\emph{not} rigid on $V(G)\setminus\{v\}$: it is missing both the short-circuit
$e_0$ and $v$'s two edges, so its rows do \emph{not} span the full
$D(|V(G)|-2)$-dimensional rigidity-row space of the rigid $G_v^{ab}$. Equivalently,
the $e_0$-row of $G_v^{ab}$ is genuinely needed in $G_v^{ab}$ --- it is \emph{not}
redundant there --- so no $e_0$-free old block of the full count exists in
$G_v^{ab}$ alone. The $e_0$-row only becomes recoverable \emph{in $G$}, after $v$'s
two new hinge rows are present; the row-counting bookkeeping that the green
sub-nodes N7b-0/1/2/3 set up cannot express this, because it asks for the old
block before $v$ is re-attached.

\medskip\noindent\emph{The corrected, motion-side rank-lift.} The realization
motive is $V(G)$-relative ($\mathrm{IsInfinitesimallyRigidOn}\,V(G)$,
\cref{def:rank-hypothesis}), so the producer is most directly a
\emph{motion-space} statement, in the established idiom of the Case~I splice
\emph{glue} \cref{lem:case-I-splice-seed} (re-adding edges only shrinks the
motion space, \cref{lem:motions-mono-of-graph-le}): \emph{if $G_v^{ab}$ is
rigid on $V(G)\setminus\{v\}$ and $v$'s two re-attaching hinges $e_a = va$,
$e_b = vb$ are transversal with $C(e_a), C(e_b)$ linearly independent, then $G$ is
rigid on $V(G)$.} The argument bypasses the $e_0$ bookkeeping entirely:
\begin{enumerate}
  \item It suffices to be rigid on $V(G)$ for the subgraph
    $G'' := (G-v) \cup \{e_a, e_b\} \le G$ (re-adding $G$'s remaining edges only
    shrinks the motion space, \cref{lem:motions-mono-of-graph-le}); $G''$ shares
    $G - v$ with $G_v^{ab}$.
  \item Let $S$ be an infinitesimal motion of $G''$ constant on
    $V(G)\setminus\{v\}$ (the hypothesis of rigid-on); in particular $S a = S b$.
    On $G - v$ this is also a motion of $G_v^{ab}$, so rigidity of $G_v^{ab}$ gives
    $S$ constant on $V(G)\setminus\{v\}$ (already assumed) --- the inductive input.
  \item The two new hinge constraints give
    $S v - S a \in \operatorname{span} C(e_a)$ and
    $S v - S b = S v - S a \in \operatorname{span} C(e_b)$, so
    $S v - S a \in \operatorname{span} C(e_a) \cap \operatorname{span} C(e_b)$.
    With $C(e_a), C(e_b)$ \emph{linearly independent}, the two $1$-dimensional spans
    meet only at $0$, so $S v = S a$: $v$ is pinned, $S$ is constant on $V(G)$.
\end{enumerate}
The general-position input ``$C(e_a), C(e_b)$ linearly independent'' is exactly the
$m = 2 \le D$ case of \cref{lem:exists-independent-panel-extensor}
($\mathrm{exists\_independent\_panelSupportExtensor}$), realized by choosing $v$'s
panel normal off the relevant locus --- the same satisfiable general-position
output N7b-1 (\cref{lem:case-II-placement-new-rows}) already consumes for the new
block. This is the genuine geometric heart of why a $1$-extension preserves rank
(KT~(6.12)), and unlike the row-side framing it is \emph{provable}: the meet step
is the elementary ``two distinct lines through the origin intersect trivially'',
not a deep extensor computation.

\medskip\noindent\emph{Buildable sub-node decomposition (N7b-4, motion-side).}
\begin{enumerate}
  \item[(M1)] \cref{lem:case-II-placement-disjoint-line-meet}: $C, C' \ne 0$ linearly
    independent $\Rightarrow \operatorname{span}\{C\} \cap \operatorname{span}\{C'\}
    = 0$ (general fact about $1$-dimensional spans; the meet step). Pure linear
    algebra, mathlib-supported.
  \item[(M2)] \cref{lem:case-II-placement-pin-vertex}: a $G''$-motion constant on
    $V(G)\setminus\{v\}$, with $v$ joined to $a, b$ by transversal hinges with
    independent support extensors, is constant on all of $V(G)$ (the pin, via M1 +
    the two hinge constraints). The crux.
  \item[(M3)] \cref{lem:case-II-placement-motion-side-assembly}: assemble M2 with the
    subgraph monotonicity \cref{lem:motions-mono-of-graph-le} and the inductive
    realization's rigidity on $V(G)\setminus\{v\}$ to conclude $G$ rigid on $V(G)$;
    feed the resulting $\mathrm{HasFullRankRealization}\,G$ to
    \cref{lem:case-II-realization}. The general-position normal for $v$ comes from
    \cref{lem:exists-independent-panel-extensor} ($m = 2$).
\end{enumerate}
This re-frame supersedes both the N7b row-side assembly path
(\cref{lem:case-II-realization-placement}) and the device-closure glue
\cref{lem:realization-of-independent-rows} \emph{for Case~II}: the motion-side
producer concludes $\mathrm{IsInfinitesimallyRigidOn}\,V(G)$ directly from the
inductive rigidity, with no detour through a witnessed independent row family.
(The green row-side sub-nodes N7b-0/1/2/3 and the glue N7a remain valid lemmas ---
they are reused by the Case~I producer, \cref{lem:case-I-realization}, which has no
single re-attached body and genuinely needs the row-counting/device route.) Build
order: M1 $\to$ M2 $\to$ M3 $\to$ N7. The motion-side route warrants its own
recon-before-build pass on M2 (the pin); M1 and M3 are mechanical.

\begin{lemma}[N7b-4 (superseded, row-side): recover $e_0$'s constraint from $v$'s two new edges; KT eq.~(6.12)]
  \label{lem:case-II-placement-e0-recovery}
  \uses{def:panel-hinge-framework, lem:case-II,
        lem:case-II-placement-old-rows-extract, lem:case-II-placement-old-rows,
        lem:case-II-placement-new-rows}
  \emph{Superseded by the motion-side re-frame above (M1--M3); kept for the audit
  trail.} This row-side formulation was found geometrically unbuildable on
  2026-06-04: it asks for an $e_0$-free independent old block of $G_v^{ab}$ of full
  count $D(|V(G)|-2)$, but no such block exists --- the $e_0$-row is not redundant
  in $G_v^{ab}$, since $G - v$ (the $e_0$-free subgraph) is not rigid (see ``Why the
  row-side framing fails'' above). The genuine rank-lift content --- recovering
  $e_0$'s span-membership constraint from $v$'s two new edges along the path
  $a$--$v$--$b$ --- is correct, but it lives \emph{in $G$} (motion-side), not as a
  row block of $G_v^{ab}$ alone, and is carried by \cref{lem:case-II-placement-pin-vertex}
  (M2). \emph{Red, superseded} (Phase~21b realization layer).
\end{lemma}
\begin{proof}
  Superseded; see the motion-side rank-lift (M1--M3) above.
\end{proof}

\begin{lemma}[M1: independent extensors have disjoint $1$-dimensional spans]
  \label{lem:case-II-placement-disjoint-line-meet}
  \uses{def:hinge-constraint}
  For $C, C' \in \mathrm{ScrewSpace}\,k$ linearly independent (in particular both
  nonzero), the $1$-dimensional spans meet trivially:
  $\operatorname{span}\{C\} \cap \operatorname{span}\{C'\} = \{0\}$. This is the meet
  step of the Case~II pin (\cref{lem:case-II-placement-pin-vertex}): a vector lying
  in both hinge spans is a scalar multiple of $C$ \emph{and} of $C'$, which forces it
  to vanish by independence. \emph{Red} (Phase~21b); pure linear algebra.
\end{lemma}
\begin{proof}
  A vector in the intersection is $\lambda C = \mu C'$; independence of $C, C'$
  forces $\lambda = \mu = 0$.
\end{proof}

\begin{lemma}[M2: two transversal hinges with independent extensors pin a re-attached body]
  \label{lem:case-II-placement-pin-vertex}
  \uses{def:hinge-constraint, def:panel-hinge-framework,
        lem:case-II-placement-disjoint-line-meet,
        lem:exists-independent-panel-extensor}
  The crux of the $1$-extension rank-lift (KT eq.~(6.12)), motion-side. Let $F$ be a
  body-hinge framework in which the body $v$ is joined to $a, b$ by hinges $e_a, e_b$
  whose support extensors $C(e_a), C(e_b)$ are linearly independent (the
  general-position condition, $m = 2$ case of
  \cref{lem:exists-independent-panel-extensor}). Then any infinitesimal motion $S$ of
  $F$ with $S a = S b$ also has $S v = S a$. Hence a motion constant on a body set
  containing $a, b$ is constant on that set together with $v$. Proof: the hinge
  constraints give $S v - S a \in \operatorname{span} C(e_a)$ and
  $S v - S b = S v - S a \in \operatorname{span} C(e_b)$, so the difference lies in
  the trivial meet (\cref{lem:case-II-placement-disjoint-line-meet}) and vanishes.
  \emph{Red} (Phase~21b); recon-before-build.
\end{lemma}
\begin{proof}
  \uses{lem:case-II-placement-disjoint-line-meet}
  From $S a = S b$ and the two hinge constraints,
  $S v - S a \in \operatorname{span} C(e_a) \cap \operatorname{span} C(e_b) = \{0\}$
  by \cref{lem:case-II-placement-disjoint-line-meet}, so $S v = S a$.
\end{proof}

\begin{lemma}[M3: Case II realization, motion-side]
  \label{lem:case-II-placement-motion-side-assembly}
  \uses{def:panel-hinge-framework, def:rank-hypothesis, lem:reduction-step,
        lem:case-II-placement-pin-vertex, lem:motions-mono-of-graph-le,
        lem:exists-independent-panel-extensor}
  The motion-side assembly producing $\mathrm{HasFullRankRealization}\,k\,G$ from the
  inductive realization of $G_v^{ab}$, superseding the row-side N7b path for Case~II.
  Choose $v$'s panel normal so $C(e_a), C(e_b)$ are linearly independent
  (\cref{lem:exists-independent-panel-extensor}, $m = 2$), placing the rest of $G$ at
  the inductive normals. The subgraph $G'' = (G-v)\cup\{e_a, e_b\} \le G$ shares
  $G - v$ with $G_v^{ab}$; on $G''$, a motion constant on $V(G)\setminus\{v\}$ (the
  inductive rigidity of $G_v^{ab}$, transported through $G - v$) is pinned at $v$ by
  \cref{lem:case-II-placement-pin-vertex}, so $G''$ is rigid on $V(G)$; re-adding $G$'s
  remaining edges keeps it rigid (\cref{lem:motions-mono-of-graph-le}). Feeds the
  resulting realization to \cref{lem:case-II-realization}. \emph{Red} (Phase~21b);
  the heart is M2.
\end{lemma}
\begin{proof}
  \uses{lem:case-II-placement-pin-vertex, lem:motions-mono-of-graph-le}
  Assemble M2 with the subgraph monotonicity
  \cref{lem:motions-mono-of-graph-le} and the inductive rigidity on
  $V(G)\setminus\{v\}$; see \cite[\S6.3]{katohTanigawa2011} (Lemma~6.8, eq.~(6.12)).
\end{proof}

